\scanpage{154}
% margin note written sideways
\margin{\textbf{Imp.\quad A. GROTHENDIECK.---Publications mathématiques, n\textsuperscript{o} 4.\\ P. U. F.\quad Placard 36}}

\textit{Proposition} \textbf{(6.4.12)}. --- \textit{Soient $f:\mathrm{X}\to \mathrm{Y}$, $g:\mathrm{Y}'\to \mathrm{Y}$ deux morphismes\,; posons $\mathrm{X}'=\mathrm{X}\times_{\mathrm{Y}}\mathrm{Y}'$ et soit $f'=f_{(\mathrm{Y}')}:\mathrm{X}'\to \mathrm{Y}'$. Soit $y'\in \mathrm{Y}'$, $y=g(y')$\,; si la fibre $f^{-1}(y)$ est un schéma algébrique fini sur $\varkappa(y)$, alors la fibre ${f'}^{-1}(y')$ est un schéma algébrique fini sur $\varkappa(y')$, ayant même rang et même nombre géométrique de points que $f^{-1}(y)$.}

Compte tenu de la transitivité des fibres (3.6.5), cela résulte aussitôt de (6.4.6) et (6.4.8).

\textbf{(6.4.13)} La prop. (6.4.11) montre que les morphismes de type fini correspondent intuitivement aux « familles algébriques de variétés algébriques », les points de $\mathrm{Y}$ jouant le rôle de « paramètres »\,; ce \deleted{sont donc les seuls} \inserted{\handwritten{qui donne à ces}} morphismes \deleted{qui aient} une signification « géométrique ». \margin{\handwritten{$\times$}} Les morphismes qui ne sont pas de type fini interviendront surtout par la suite dans les questions de « changement du préschéma de base », par exemple par localisation ou complétion.

\textbf{6.5. Détermination locale d’un morphisme.}

\textit{Proposition} \textbf{(6.5.1)}. --- \textit{Soient $\mathrm{X}$, $\mathrm{Y}$ deux $\mathrm{S}$-préschémas, $\mathrm{Y}$ étant de type fini sur $\mathrm{S}$\,; soient $x\in \mathrm{X}$, $y\in \mathrm{Y}$ au-dessus d’un même point $s\in \mathrm{S}$.}

(i) \textit{Si deux $\mathrm{S}$-morphismes $f=(\psi,\theta)$, $f'=(\psi',\theta')$ de $\mathrm{X}$ dans $\mathrm{Y}$ sont tels que $\psi(x)=\psi'(x)=y$, et que les $\mathcal{O}_s$-homomorphismes (locaux) $\theta_x^{\sharp}$ et ${\theta'}^{\sharp}_x$ de $\mathcal{O}_y$ dans $\mathcal{O}_x$ soient identiques, $f$ et $f'$ coïncident dans un voisinage ouvert de $x$.}

(ii) \textit{Supposons en outre $\mathrm{S}$ localement noethérien. Pour tout $\mathcal{O}_s$-homomorphisme local $\varphi:\mathcal{O}_y\to\mathcal{O}_x$, il existe un voisinage ouvert $\mathrm{U}$ de $x$ dans $\mathrm{X}$ et un $\mathrm{S}$-morphisme $f=(\psi,\theta)$ de $\mathrm{U}$ dans $\mathrm{Y}$ tels que $\psi(x)=y$ et $\theta_x^{\sharp}=\varphi$.}

(i) La question étant locale sur $\mathrm{S}$, $\mathrm{X}$ et $\mathrm{Y}$, on peut supposer $\mathrm{S}$, $\mathrm{X}$, $\mathrm{Y}$ affines d’anneaux respectifs $\mathrm{A}$, $\mathrm{B}$, $\mathrm{C}$, $f$ et $f'$ étant de la forme $({}^{a}\varphi,\widetilde{\varphi})$ et $({}^{a}\varphi',\widetilde{\varphi}')$ respectivement, où $\varphi$ et $\varphi'$ sont deux $\mathrm{A}$-homomorphismes de $\mathrm{C}$ dans $\mathrm{B}$ tels que $\varphi^{-1}(\unsure{\mathfrak{j}|\mathrm{j}}_x)={\varphi'}^{-1}(\mathfrak{j}_x)=\mathfrak{j}_y$, et les homomorphismes $\varphi_x$ et $\varphi'_x$ de $\mathrm{C}_y$ dans $\mathrm{B}_x$, déduits de $\varphi$ et $\varphi'$, sont identiques\,; on peut en outre supposer que $\mathrm{C}$ est une $\mathrm{A}$-algèbre \textit{de type fini}. Soient $c_i$ $(1\leqslant i\leqslant n)$ des générateurs de la $\mathrm{A}$-algèbre $\mathrm{C}$, et posons $b_i=\varphi(c_i)$, $b'_i=\varphi'(c_i)$\,; par hypothèse, on a $b_i/1=b'_i/1$ dans l’anneau de fractions $\mathrm{B}_x$ $(1\leqslant i\leqslant n)$. Cela signifie qu’il existe des éléments $s_i\in \mathrm{B}-\mathfrak{j}_x$ tels que $s_i(b_i-b'_i)=0$ pour $1\leqslant i\leqslant n$, et on peut évidemment supposer tous les $s_i$ égaux à un même élément $g\in \mathrm{B}-\mathfrak{j}_x$. On en conclut que l’on a $b_i/1=b'_i/1$ pour $1\leqslant i\leqslant n$ dans l’anneau de fractions $\mathrm{B}_g$\,; si $i_g$ est l’homomorphisme canonique $\mathrm{B}\to \mathrm{B}_g$, on a par suite $i_g\circ\varphi=i_g\circ\varphi'$\,; donc les restrictions de $f$ et $f'$ à $\mathrm{D}(g)$ sont identiques.

(ii) On peut se ramener à la même situation que dans (i), et supposer en outre que l’anneau $\mathrm{A}$ est \textit{noethérien}. Soient $c_i$ $(1\leqslant i\leqslant n)$ des générateurs de la $\mathrm{A}$-algèbre $\mathrm{C}$, et soit $\alpha:\mathrm{A}[\mathrm{X}_1,\ldots,\mathrm{X}_{\unsure{n|i}}]\to \mathrm{C}$ l’homomorphisme de l’algèbre de polynômes $\mathrm{A}[\mathrm{X}_1,\ldots,\mathrm{X}_n]$ sur $\mathrm{C}$ transformant $\mathrm{X}_i$ en $c_i$ pour $1\leqslant i\leqslant n$. Soit d’autre part $i_y$ l’homomorphisme canonique $\mathrm{C}\to \mathrm{C}_y$, et considérons l’homomorphisme composé
\[
\beta:\mathrm{A}[\mathrm{X}_1,\ldots,\mathrm{X}_n]\xrightarrow{\alpha}\mathrm{C}\xrightarrow{\unsure{iy|i_y}}\mathrm{C}_y\xrightarrow{\varphi}\mathrm{B}_x.
\]
% archive sheet number: 153
